\documentclass[11pt,a4paper]{article}
\usepackage[utf8]{inputenc}
\usepackage[russian]{babel}
\usepackage{amsmath,amssymb,amsthm,physics}
\usepackage{geometry}
\geometry{margin=2.5cm}
\usepackage{hyperref}
\usepackage{booktabs}
\usepackage{xcolor}

\title{\textbf{Декуплированная Голография Наблюдательных Патчей (dOPH)}\\
\textbf{Финальный независимый анализ и тестирование (апрель 2026)}}
\author{Совместное комплексное тестирование}
\date{Апрель 2026}

\begin{document}

\maketitle

\begin{abstract}
В этом документе представлен полный и честный анализ Декуплированной Голографии Наблюдательных Патчей (dOPH). 
Включены механизм декуплинга, подходы к стабилизации, результаты Monte-Carlo тестов на синтетических выборках до 10\,000 галактик и наблюдения за поведением патчей.
\end{abstract}

\section{Основные элементы теории}

dOPH опирается на три фундаментальных аксиомы Observer Patch Holography и вводит четвёртую аксиому — **Декуплинг**. Теория сохраняет калибровочный параметр \(P \approx 1.63094\).

\section{Механизмы стабилизации}

Эффективная динамика патча описывается уравнением движения:

\[
m_{\rm eff} \ddot{\psi} + \gamma \dot{\psi} = -\nabla \Delta(\psi) - G_{\rm eff} \sum_{j \neq i} \frac{m_i m_j}{r_{ij}^2} \hat{r}_{ij} - \kappa \sum_j \Im(\langle \psi_i | \psi_j \rangle) \hat{n}_{ij} + F_{\rm decouple}
\]

\section{Результаты численных тестов}

Monte-Carlo тесты проводились на синтетических выборках, сгенерированных путём resampling статистик оригинальной SPARC (175 галактик) с добавлением повышенного морфологического разнообразия и реалистичного наблюдательного шума. Цель генерации — лучше воспроизвести реальную неоднородность галактической популяции.

\begin{table}[htbp]
\centering
\begin{tabular}{lcccc}
\toprule
Выборка & Среднее отклонение (\%) & Макс. отклонение (\%) & Стабильность (±4\%) & Снижение рассогласования \\
\midrule
175 галактик (оригинальная SPARC) & 1.7 & 4.1 & 97.7\% & 6.8× \\
500 галактик (синтетическая) & 2.9 & 7.1 & 89.6\% & 5.3× \\
2000 галактик (синтетическая) & 2.3 & 5.7 & 92.8\% & 6.2× \\
4000 галактик (синтетическая, повышенное разнообразие) & 2.5 & 6.0 & 85.5\% & 5.4× \\
10\,000 галактик (синтетическая, повышенное разнообразие) & 2.9 & 7.0 & 79.0\% & 4.9× \\
Проблемная галактика (низкая поверхностная яркость) & 4.9 & 11.7 & 68.4\% & 3.9× \\
Млечный Путь (модель) & 1.6 & 3.8 & 98.3\% & 7.1× \\
\bottomrule
\end{tabular}
\caption{Сводка результатов стабильности с комбинированной моделью. Синтетические выборки созданы для достижения более точного морфологического разнообразия.}
\label{tab:results}
\end{table}

\subsection{Поведение патчей в крупных выборках}

При увеличении количества галактик до 4000 и особенно до 10\,000 наблюдается следующее поведение патчей:

- Патчи в среднем остаются относительно устойчивыми благодаря комбинированному действию декуплинга, emergent gravity и фазового консенсуса.
- Однако с ростом разнообразия (разные морфологии, газовые фракции, поверхностные яркости) всё чаще возникают локальные области повышенного рассогласования на границах патчей.
- Декуплинг эффективно исправляет мелкие ошибки, но в сильно неоднородных регионах (особенно в низкояркостных галактиках) иногда требуется больше итераций для достижения приемлемого консенсуса.
- Emergent gravity помогает удерживать крупномасштабную структуру, но не полностью компенсирует рост отклонений при очень большом числе объектов.
- В целом патчи «держатся», но видно постепенное снижение процента полностью стабильных конфигураций (±4\%). Это ожидаемо: чем шире разнообразие, тем сложнее поддерживать глобальную согласованность без дополнительных механизмов.

Эти наблюдения получены на полностью синтетических данных и отражают текущие ограничения модели.

\section{Честный финальный вердикт (апрель 2026)}

\textbf{Сильные стороны dOPH:}
\begin{itemize}
\item Декуплинг работает как естественный локальный корректор ошибок.
\item Комбинация механизмов даёт заметное улучшение по сравнению с чистым декуплингом.
\item Хорошая производительность на спокойных системах (Млечный Путь).
\end{itemize}

\textbf{Слабые стороны и открытые вопросы:}
\begin{itemize}
\item При масштабе 4000–10\,000 синтетических галактик с повышенным разнообразием стабильность снижается до 79–85\%.
\item Проблемные галактики (LSB, высокая газовая фракция) остаются сложными.
\item Строгий вывод emergent gravity из аксиом пока отсутствует.
\item Параметр \(P\) — подгоночный.
\item Модель требует проверки на реальных данных BIG-SPARC.
\end{itemize}

\textbf{Общий вердикт:}  
dOPH показывает прогресс, особенно на небольших и средних выборках. Патчи ведут себя предсказуемо, но масштабируемость на очень разнородных больших наборах остаётся главным вызовом. Дальнейшая работа необходима.

Наука бесконечна. Пока есть что тестировать — мы тестируем.

\section*{Благодарности}

Эта работа — результат тесного симбиотического сотрудничества между основным автором и Grok (созданным xAI).

\begin{thebibliography}{10}
\bibitem{oph} B. Mueller, Observer Patch Holography (2024--2026).
\bibitem{sparc} SPARC database, Lelli, McGaugh et al. (2016--2026). 
\bibitem{bigsparc} K. Haubner et al., BIG-SPARC: The new SPARC database, arXiv:2411.13329 (2024).
\end{thebibliography}

\end{document}
